% !TeX encoding = UTF-8
\documentclass[variant=docente,cover=institutional,mode=screen]{ua-engineering-v6-article}
\usepackage{longtable,array}
% !TeX encoding = UTF-8
\UASetUniversity{Universidad Austral}
\UASetFaculty{Facultad de Ingeniería}
\UASetDepartment{Departamento de Computación}
\UASetCampus{Pilar}
\UASetCareer{Ingeniería Informática}
\UASetCommission{Comisión A}
\UASetProfessor{Equipo docente}
\UASetSemester{Segundo cuatrimestre 2026}
\UASetCourse{Sistemas Distribuidos}
\UASetDate{2026}

\UASetClassNumber{08}
\UASetClassTitle{Causalidad, consenso y consistencia distribuida}
\title{Clase 08 --- Guía docente y libreto de exposición}
\author{\UAProfessor}
\date{\UADate}
\newcommand{\FrameField}[2]{\par\noindent\textbf{#1.} #2\par}
\begin{document}
\maketitle

\section{Introducción narrativa}
En esta clase seguiremos el destino de una única entrada VIP que Sofía y Tomás intentan comprar casi al mismo tiempo desde regiones diferentes. Como los nodos no comparten una visión instantánea del sistema ni pueden confiar en un reloj global perfecto, primero utilizaremos causalidad y relojes lógicos para comprender qué eventos están relacionados y cuáles son concurrentes. Luego veremos cómo Raft y Paxos permiten que varias réplicas acuerden una sola decisión autoritativa, y cómo CAP y PACELC obligan a decidir entre esperar, rechazar, responder con información potencialmente atrasada o reparar más tarde. Finalmente, compararemos dos estrategias reales: Dynamo, que favorece disponibilidad conservando versiones y pagando reconciliación, y Spanner, que utiliza consenso y TrueTime para sostener transacciones globales con un orden compatible con lo observado por los usuarios. La pregunta que une toda la clase será: qué error resulta inaceptable para cada operación y qué costo ---latencia, indisponibilidad, coordinación o reparación--- estamos dispuestos a pagar para evitarlo.

\section{Objetivo pedagógico profundo}
La clase debe transformar siglas y algoritmos en decisiones sobre una ejecución concreta. El estudiante debe terminar pudiendo explicar qué sabe cada nodo, qué historia incorrecta se evita, qué condición habilita progreso y qué costo operacional sostiene la garantía.

\section{Objetivos de la clase}
Al finalizar, el estudiante debe poder:
\begin{itemize}
\item explicar la diferencia entre preservación de registros y decisión autoritativa;
\item reconstruir \emph{happened-before}, Lamport y relojes vectoriales;
\item explicar máquina de estados, safety, liveness y la secuencia de commit de Raft;
\item transferir el problema de consenso al vocabulario de Paxos/Multi-Paxos;
\item distinguir log de consenso, WAL local y log de eventos;
\item derivar CAP a partir de historias indistinguibles y explicar PACELC;
\item diferenciar recencia, linearizabilidad, serializabilidad y consistencia externa;
\item separar mayoría de consenso de quórums Dynamo-like;
\item razonar sobre versiones concurrentes, sloppy quorum y reparación;
\item explicar TrueTime, timestamp de commit, commit wait y orden externo;
\item defender una política con invariante, falla, costo y métrica.
\end{itemize}

\subsection{Resultados de aprendizaje observables}
Al finalizar, el estudiante deberá poder reconstruir una ejecución causal, actualizar relojes de Lamport y vectoriales, distinguir propuesta de commit en Raft, explicar el trade-off CAP para una operación concreta, comparar mayoría de consenso con quórum Dynamo-like y justificar el commit wait de Spanner. La evidencia será un timeline o defensa que declare invariante, falla, frontera de respuesta, costo y métrica.

\section{Parte A: guion operativo}
Esta parte organiza la conducción de las tres horas: apertura, preparación, ruta núcleo, guion por diapositiva, demostraciones, preguntas y contingencias. El objetivo es que la progresión vaya de la intuición visual a la formalización sin perder el caso VIP-01.

\section{Arquitectura didáctica v9 para 180 minutos}
\begin{longtable}{p{0.13\textwidth}p{0.23\textwidth}p{0.56\textwidth}}
\toprule
Tiempo & Bloque & Propósito\\
\midrule
0--12 & Apertura/Kafka & Recuperar logs, particiones, replay y durabilidad; formular el límite de Kafka.\\
12--27 & Relevancia/modelo & etcd, Dynamo, Spanner; réplica, aislamiento, caída y timeout.\\
27--55 & Causalidad & Happened-before, Lamport, vectores y tres tipos de reloj.\\
55--90 & Consenso & Máquina de estados, Raft, crash, total order, GC/fencing y Paxos.\\
90--100 & Pausa & Reconstrucción por parejas del log 87 y preguntas anónimas.\\
100--122 & CAP/PACELC & Partición 3|2, historias indistinguibles, disponibilidad y garantías.\\
122--150 & Dynamo & Quórums, contraejemplos, versiones y reparación.\\
150--168 & Spanner & Grupos, TrueTime, commit wait y consistencia externa.\\
168--177 & Integración & Cuatro operaciones, cuatro políticas y pruebas.\\
177--180 & Exit ticket & Evidencia individual y dudas para Clase 09.\\
\bottomrule
\end{longtable}

\begin{uawarn}
La fuente de diapositivas es invariante. No elimine láminas para acortar: use una ruta núcleo y deje las láminas de extensión/respaldos para preguntas. Los tiempos del libreto indican qué slides deben recorrerse rápidamente.
\end{uawarn}

\section{Preparación y plan de pizarrón}
\subsection{Preparación previa}
\begin{itemize}
\item Izquierda: tres líneas BA--Líder--Madrid; causalidad, Lamport y vectores.
\item Centro: N1--N5, entry 87, term 12, commitIndex y posición del crash.
\item Derecha: invariante, respuesta durante partición, costo normal y métrica.
\item Llevar tarjetas N1--N5, SOFÍA, TOMÁS, APPEND, ACK, COMMIT, CRASH y PARTICIÓN.
\item Ejecutar previamente las demos de causalidad, Raft, quórums y TrueTime.
\item Mantener visible durante toda la clase: \(confirmaciones(VIP\mbox{-}01)\le 1\).
\end{itemize}

\section{Guion de apertura}
Diga: ``Hoy no estudiaremos algoritmos aislados. Seguiremos una única entrada VIP que Sofía y Tomás intentan comprar desde regiones diferentes. Primero describiremos qué eventos pueden haberse influido; después veremos cómo cinco réplicas eligen una sola decisión; finalmente compararemos diseños que esperan, rechazan, conservan versiones o reparan.''

Pida una predicción de bajo riesgo: \emph{si la red divide tres réplicas de dos, ¿qué debería responder el lado de dos cuando todavía ve stock igual a uno?} Recoja tres opciones sin corregir. Escriba el invariante y anuncie que cada bloque volverá a esa misma pregunta.

\subsection{La intuición que organiza la clase}
La intuición central es que un nodo no ve el mundo completo: solo posee estado local, mensajes recibidos y timeouts. Los relojes lógicos describen conocimiento; el consenso elige autoridad; CAP muestra la decisión durante una partición; Dynamo conserva conflictos reparables; Spanner expone incertidumbre temporal y espera cuando necesita orden externo.

\section{Ruta núcleo y láminas de respaldo}
La ruta núcleo incluye apertura, Kafka, caso VIP, causalidad, Lamport, vectores, máquina de estados, Raft, CAP, diferencia de quórums, un contraejemplo, version vectors, flujo Dynamo, TrueTime, integración y exit ticket. Paxos detallado, Kubernetes, tres contraejemplos completos y mecanismos de repair pueden acelerarse o usarse como respaldo según el ritmo.

\section{Guion por diapositiva}

\section{Bloque: Apertura}

\subsection*{C08-00 --- La historia que guiará toda la clase}
\FrameField{Propósito}{Instalar el conflicto único y la pregunta de costos.}
\FrameField{Qué decir}{Lea el primer párrafo casi literalmente. Subraye que la clase no es una lista de productos: seguiremos el mismo recurso desde mensajes concurrentes hasta transacciones globales.}
\FrameField{Pregunta}{¿Qué sería peor para una entrada única: esperar o confirmar dos dueños?}
\FrameField{Respuesta esperada}{Depende del dominio, pero dos dueños viola el invariante; esperar/rechazar puede ser aceptable.}
\FrameField{Cierre/transición}{Antes de estudiar mecanismos, explicaremos qué evidencias deben producir.}
\FrameField{Tiempo sugerido}{2 min}

\subsection*{C08-00B --- Evidencias de aprendizaje}
\FrameField{Propósito}{Hacer observables los resultados y reducir ansiedad.}
\FrameField{Qué decir}{Explique que no se evaluará recitar CAP o Raft, sino reconstruir ejecuciones: flechas causales, clocks, estado de logs, mayoría y comportamiento bajo falla.}
\FrameField{Pregunta}{¿Qué producción concreta demostraría que entendieron Raft?}
\FrameField{Respuesta esperada}{Poder mover un crash y justificar si la entrada sobrevive, se aplica o se responde.}
\FrameField{Cierre/transición}{Comenzamos recuperando Kafka para separar preservar eventos de decidir autoridad.}
\FrameField{Tiempo sugerido}{2 min}

\section{Bloque: Puente desde Kafka}

\subsection*{C08-KF01 --- Repaso visual: un Mundial en tiempo real}
\FrameField{Propósito}{Reactivar productor, log y consumidor mediante un caso simple.}
\FrameField{Qué decir}{Recorra de izquierda a derecha. Defina evento como hecho, registro como su representación y log como secuencia preservada y releíble.}
\FrameField{Pregunta}{¿Por qué el marcador no necesita bloquear al productor del partido?}
\FrameField{Respuesta esperada}{Porque el registro queda preservado y el consumidor puede avanzar/reintentar independientemente.}
\FrameField{Cierre/transición}{Ahora aumentaremos la carga para revelar la necesidad de particionar.}
\FrameField{Tiempo sugerido}{1.5 min}

\subsection*{C08-KF02 --- Una única secuencia no alcanza para muchos partidos}
\FrameField{Propósito}{Motivar particionamiento por capacidad y orden.}
\FrameField{Qué decir}{Separe dos problemas: una sola cola concentra recursos; repartir al azar puede romper el orden por partido. No ofrezca todavía la solución.}
\FrameField{Pregunta}{¿Qué unidad necesita conservar orden: todo el Mundial o cada partido?}
\FrameField{Respuesta esperada}{Cada partido; partidos distintos pueden avanzar en paralelo.}
\FrameField{Cierre/transición}{La key representará ese ámbito de orden.}
\FrameField{Tiempo sugerido}{1.5 min}

\subsection*{C08-KF03 --- La key define el ámbito de orden}
\FrameField{Propósito}{Recordar que key es una decisión semántica, no solo hash.}
\FrameField{Qué decir}{Muestre P0, P1 y P2 como logs independientes. \texttt{game\_id} mantiene juntas las transiciones de un partido, pero no crea orden global entre partidos.}
\FrameField{Pregunta}{¿Qué key produciría una hot partition?}
\FrameField{Respuesta esperada}{Una key demasiado gruesa o muy frecuente, por ejemplo \texttt{tournament\_id} para todo el Mundial.}
\FrameField{Cierre/transición}{Una vez particionado, el trabajo puede repartirse entre consumidores.}
\FrameField{Tiempo sugerido}{2 min}

\subsection*{C08-KF04 --- Consumer groups: repartir trabajo y releer la misma historia}
\FrameField{Propósito}{Recuperar reparto dentro de grupo y fan-out entre grupos.}
\FrameField{Qué decir}{Dibuje mentalmente dos niveles: dentro de live-score cada partición tiene un único dueño activo; analytics usa otro group.id y relee toda la historia con offsets propios.}
\FrameField{Pregunta}{¿Qué pasa si hay cinco consumidores y tres particiones?}
\FrameField{Respuesta esperada}{Dos quedan ociosos; el paralelismo activo está acotado por las particiones.}
\FrameField{Cierre/transición}{La recuperación del consumidor revela duplicados y replay.}
\FrameField{Tiempo sugerido}{2 min}

\subsection*{C08-KF05 --- Crash, replay e idempotencia}
\FrameField{Propósito}{Reactivar el timeline poll--efecto--commit.}
\FrameField{Qué decir}{Mueva el crash después del efecto y antes del committed offset. El replay es correcto para no perder, pero puede repetir el efecto. Introduzca \texttt{event\_id}/inbox.}
\FrameField{Pregunta}{¿Qué sería durable para impedir emitir dos QR?}
\FrameField{Respuesta esperada}{La identidad procesada y el efecto deben registrarse atómicamente o el efecto debe ser naturalmente idempotente.}
\FrameField{Cierre/transición}{Antes de abandonar Kafka, ubicaremos números y configuración de durabilidad.}
\FrameField{Tiempo sugerido}{2 min}

\subsection*{C08-KF06 --- Algunas métricas a considerar en Kafka}
\FrameField{Propósito}{Aclarar RF/ISR/acks y separar guardrails de límites universales.}
\FrameField{Qué decir}{Explique RF=3, ISR, minISR=2 y acks=all con una réplica fuera. Insista: si queda una sola en ISR, se rechaza. Los valores de 1 MB, 1 TB y 10 MB/s son presupuestos docentes, no máximos de Kafka.}
\FrameField{Pregunta}{¿Por qué acks=all no significa necesariamente tres ACK?}
\FrameField{Respuesta esperada}{Porque espera a todas las réplicas del ISR actual, siempre que su tamaño cumpla minISR.}
\FrameField{Cierre/transición}{Kafka protege el registro; todavía falta elegir una única decisión de negocio.}
\FrameField{Tiempo sugerido}{3 min}

\subsection*{C08-KF07 --- Puente: Kafka ordena registros; no elige por sí solo una decisión autoritativa}
\FrameField{Propósito}{Cerrar el repaso y abrir consenso.}
\FrameField{Qué decir}{Contraste dos preguntas: Kafka preserva OrderAttempt(Sofía) y OrderAttempt(Tomás); el servicio de reservas debe decidir cuál cambia AVAILABLE. Dos particiones independientes no pueden confirmar dos propietarios.}
\FrameField{Pregunta}{¿Podemos resolver la última entrada usando solo una buena key?}
\FrameField{Respuesta esperada}{Solo si existe una autoridad/serialización clara detrás; la key por sí sola no crea consenso entre réplicas.}
\FrameField{Cierre/transición}{Veamos sistemas actuales que pagan estos costos en producción.}
\FrameField{Tiempo sugerido}{1.5 min}

\section{Bloque: Por qué estos conceptos se usan hoy}

\subsection*{C08-R01 --- Sistema real 1: Kubernetes y etcd}
\FrameField{Propósito}{Definir producto y ubicar consenso en un uso actual.}
\FrameField{Qué decir}{Defina Kubernetes como orquestador y etcd como almacén autoritativo del control plane. Raft sostiene un log común para cambios de configuración; sin mayoría no se confirman cambios nuevos.}
\FrameField{Pregunta}{¿Qué puede seguir ocurriendo si etcd pierde quorum?}
\FrameField{Respuesta esperada}{Cargas existentes pueden continuar un tiempo, pero el control plane no puede comprometer nuevo estado con normalidad.}
\FrameField{Cierre/transición}{Dynamo ejemplifica otra prioridad: seguir respondiendo y reparar.}
\FrameField{Tiempo sugerido}{2 min}

\subsection*{C08-R02 --- Sistema real 2: Amazon Dynamo}
\FrameField{Propósito}{Definir Dynamo antes de sus mecanismos.}
\FrameField{Qué decir}{Describa el almacén clave--valor altamente disponible del paper, no DynamoDB actual. El carrito permite conservar ramas concurrentes porque el merge es semánticamente posible.}
\FrameField{Pregunta}{¿Por qué no aplicar la misma unión a VIP-01?}
\FrameField{Respuesta esperada}{Porque unir dos propietarios viola un invariante no fusionable.}
\FrameField{Cierre/transición}{Spanner mostrará una estrategia que paga más coordinación por orden fuerte global.}
\FrameField{Tiempo sugerido}{2 min}

\subsection*{C08-R03 --- Sistema real 3: Google Spanner}
\FrameField{Propósito}{Definir Spanner y anticipar el problema del tiempo global.}
\FrameField{Qué decir}{Diga que es una base distribuida global con consenso por grupos y transacciones multi-grupo. TrueTime no es un reloj perfecto: expone incertidumbre para que el sistema espere cuando necesita orden externo.}
\FrameField{Pregunta}{¿Por qué un timestamp de pared común no bastaría?}
\FrameField{Respuesta esperada}{Porque los relojes tienen error/deriva y pueden violar el orden real si se tratan como exactos.}
\FrameField{Cierre/transición}{Ahora convertimos estos ejemplos en decisiones de ingeniería.}
\FrameField{Tiempo sugerido}{2 min}

\subsection*{C08-R04 --- Lo que un ingeniero decide realmente}
\FrameField{Propósito}{Orientar toda la clase hacia operaciones e invariantes.}
\FrameField{Qué decir}{Recorra la matriz: última entrada, carrito, ranking, ledger. Muestre que la misma infraestructura puede ofrecer políticas distintas por dato y operación.}
\FrameField{Pregunta}{¿Qué error tolera ranking que no tolera un ledger?}
\FrameField{Respuesta esperada}{Atraso temporal; el ledger no tolera crear o perder dinero.}
\FrameField{Cierre/transición}{El mapa de la clase ordena los mecanismos alrededor de esas decisiones.}
\FrameField{Tiempo sugerido}{2 min}

\subsection*{C08-MAP --- Mapa de la clase}
\FrameField{Propósito}{Anunciar la progresión y evitar sensación de temas aislados.}
\FrameField{Qué decir}{Señale el recorrido: observación parcial, causalidad, consenso, partición, versiones y orden global. Vuelva a este mapa en cada transición.}
\FrameField{Pregunta}{¿Qué bloque elige una decisión y cuál detecta ramas concurrentes?}
\FrameField{Respuesta esperada}{Raft/Paxos eligen; vector clocks detectan causalidad/concurrencia.}
\FrameField{Cierre/transición}{Comenzamos con el modelo mínimo de una réplica.}
\FrameField{Tiempo sugerido}{1 min}

\section{Bloque: Modelo mínimo y caso conductor}

\subsection*{C08-M01 --- Modelo mínimo: qué contiene una réplica}
\FrameField{Propósito}{Compensar brechas de SO/hardware con un modelo suficiente.}
\FrameField{Qué decir}{Distinga proceso, RAM, medio persistente y red. No profundice en MMU. Pregunte qué se pierde en cada tipo de falla.}
\FrameField{Pregunta}{¿Qué sobrevive a un crash del proceso y qué a un corte de energía?}
\FrameField{Respuesta esperada}{Depende de la frontera persistente; RAM no, medio persistente correctamente sincronizado sí.}
\FrameField{Cierre/transición}{Con ese modelo separaremos aislamiento de caída.}
\FrameField{Tiempo sugerido}{2 min}

\subsection*{C08-M02 --- Tres situaciones que no debemos confundir}
\FrameField{Propósito}{Diferenciar sano, aislado y caído.}
\FrameField{Qué decir}{Explique que un nodo aislado puede ejecutar y modificar estado local, lo que lo hace más peligroso que un nodo claramente muerto. Un timeout no distingue estas situaciones.}
\FrameField{Pregunta}{¿Cuál situación puede producir dos autoridades activas?}
\FrameField{Respuesta esperada}{Una partición o pausa prolongada seguida de elección, si no existen términos/fencing correctos.}
\FrameField{Cierre/transición}{El timeout es nuestra primera observación ambigua.}
\FrameField{Tiempo sugerido}{2 min}

\subsection*{C08-M03 --- Timeout: observación local, no diagnóstico}
\FrameField{Propósito}{Instalar indistinguibilidad como fundamento.}
\FrameField{Qué decir}{Recorra las cuatro historias. No permita que digan “el servidor cayó” sin evidencia. Conecte con idempotencia y reintentos.}
\FrameField{Pregunta}{¿Qué puede afirmar A con certeza?}
\FrameField{Respuesta esperada}{Solo que no recibió una respuesta válida dentro de su plazo.}
\FrameField{Cierre/transición}{Ahora aplicaremos el modelo al recurso VIP-01.}
\FrameField{Tiempo sugerido}{2 min}

\subsection*{C08-M04 --- Caso conductor: Festival Austral Global}
\FrameField{Propósito}{Fijar actores, recurso e invariante permanente.}
\FrameField{Qué decir}{Mantenga visible stock=1 y confirmaciones<=1. Aclare que “estado replicado” no implica que las copias estén siempre sincronizadas.}
\FrameField{Pregunta}{¿Qué condición debe ser verdadera después de cualquier falla?}
\FrameField{Respuesta esperada}{Como máximo una confirmación autoritativa para VIP-01.}
\FrameField{Cierre/transición}{Pediremos una predicción antes de introducir teoría.}
\FrameField{Tiempo sugerido}{2 min}

\subsection*{C08-M05 --- Predicción inicial}
\FrameField{Propósito}{Obtener compromiso cognitivo con participación de bajo riesgo.}
\FrameField{Qué decir}{Pida voto individual entre confirmar, rechazar/esperar o aceptar intención. No corrija de inmediato; solicite una frase sobre el error aceptado.}
\FrameField{Pregunta}{¿Quién paga cada opción?}
\FrameField{Respuesta esperada}{Confirmar ambos paga conflicto; rechazar paga indisponibilidad; pendiente cambia la promesa al usuario.}
\FrameField{Cierre/transición}{Para razonar mejor necesitamos describir relaciones entre eventos.}
\FrameField{Tiempo sugerido}{3 min}

\section{Bloque: Causalidad y relojes lógicos}

\subsection*{C08-C01 --- Happened-before: tres reglas}
\FrameField{Propósito}{Formalizar causalidad después del caso.}
\FrameField{Qué decir}{Defina orden local, mensaje y transitividad. Use flechas; evite hablar de tiempo físico exacto.}
\FrameField{Pregunta}{¿Dos eventos sin cadena de mensajes son simultáneos?}
\FrameField{Respuesta esperada}{No necesariamente; son concurrentes en el sentido causal.}
\FrameField{Cierre/transición}{Veamos ambas reservas y al líder en un único timeline.}
\FrameField{Tiempo sugerido}{2 min}

\subsection*{C08-C02 --- Las dos reservas nacen en ramas independientes y llegan al líder}
\FrameField{Propósito}{Visualizar la causalidad completa con tres líneas.}
\FrameField{Qué decir}{Recorra BA, líder y Madrid. El líder ordena sus recepciones, pero ese orden no vuelve causal una solicitud respecto de la otra.}
\FrameField{Pregunta}{¿Qué pares son causales y cuál es concurrente?}
\FrameField{Respuesta esperada}{a1→a2→l1 y m1→m2→l3; a2||m2.}
\FrameField{Cierre/transición}{Lamport convertirá las flechas causales en números compatibles.}
\FrameField{Tiempo sugerido}{3 min}

\subsection*{C08-C03 --- Lamport: convertir relaciones causales en números crecientes}
\FrameField{Propósito}{Explicar variables y regla de recepción.}
\FrameField{Qué decir}{Defina L y tm. Explique max como unión de pasados conocidos y +1 como nuevo evento local posterior.}
\FrameField{Pregunta}{Con L=4 y tm=7, ¿por qué no alcanza asignar 7?}
\FrameField{Respuesta esperada}{La recepción debe ser estrictamente posterior, por eso 8.}
\FrameField{Cierre/transición}{Apliquemos la regla a la misma ejecución.}
\FrameField{Tiempo sugerido}{2 min}

\subsection*{C08-X01 --- La misma ejecución con relojes de Lamport}
\FrameField{Propósito}{Aplicar Lamport sin cambiar el diagrama.}
\FrameField{Qué decir}{Calcule BA 1,2; Madrid 1,2; líder 3,4,5. Destaque que ambos envíos pueden tener 2 y seguir siendo eventos diferentes concurrentes.}
\FrameField{Pregunta}{¿Qué nos demuestra L=5 en el líder?}
\FrameField{Respuesta esperada}{Que su evento ocurre después de los pasados conocidos, no quién causó a quién entre las ramas originales.}
\FrameField{Cierre/transición}{La siguiente lámina abstrae el cálculo y su uso como desempate.}
\FrameField{Tiempo sugerido}{3 min}

\subsection*{C08-C04 --- Ejemplo Lamport sobre las dos reservas}
\FrameField{Propósito}{Consolidar el cálculo y mostrar orden total convencional.}
\FrameField{Qué decir}{Muestre que \texttt{(timestamp,process\_id)} puede desempatar para ordenar logs o trazas, pero esa convención no crea causalidad de negocio.}
\FrameField{Pregunta}{¿Un desempate determina quién merece la entrada?}
\FrameField{Respuesta esperada}{No; solo produce un orden convencional. La autoridad la decide un protocolo/estado de dominio.}
\FrameField{Cierre/transición}{Necesitamos entender el límite formal de Lamport.}
\FrameField{Tiempo sugerido}{2 min}

\subsection*{C08-C05 --- Límite de Lamport}
\FrameField{Propósito}{Distinguir implicación de equivalencia.}
\FrameField{Qué decir}{Escriba las dos fórmulas. Construya un contraejemplo con procesos sin mensajes y marcas 1 y 2.}
\FrameField{Pregunta}{¿Qué información falta en un escalar?}
\FrameField{Respuesta esperada}{Qué actor contribuyó a la historia y si dos avances son independientes.}
\FrameField{Cierre/transición}{Los vectores conservarán conocimiento por actor.}
\FrameField{Tiempo sugerido}{2 min}

\subsection*{C08-C06 --- Relojes vectoriales: registrar conocimiento por actor}
\FrameField{Propósito}{Introducir componentes y comparación parcial.}
\FrameField{Qué decir}{Explique incremento local, máximo por componente al recibir y dominación. Mantenga fijo el orden [BA,Madrid,Líder].}
\FrameField{Pregunta}{¿Qué significa que ningún vector domine?}
\FrameField{Respuesta esperada}{Que existen avances independientes: las versiones/eventos son concurrentes.}
\FrameField{Cierre/transición}{Volvamos al mismo timeline para calcularlos.}
\FrameField{Tiempo sugerido}{2 min}

\subsection*{C08-X02 --- La misma ejecución con relojes vectoriales}
\FrameField{Propósito}{Leer visualmente vectores por evento.}
\FrameField{Qué decir}{Señale [2,0,0] y [0,2,0] como incomparables. Luego siga al líder hasta [2,2,3]. No mezcle el 5 de Lamport.}
\FrameField{Pregunta}{¿Qué conocimiento nuevo aparece en l3?}
\FrameField{Respuesta esperada}{El líder ya conoce dos eventos de BA, dos de Madrid y ha ejecutado tres propios.}
\FrameField{Cierre/transición}{La siguiente lámina descompone el cálculo del líder.}
\FrameField{Tiempo sugerido}{3 min}

\subsection*{C08-X03 --- Cómo actualiza su reloj vectorial el líder}
\FrameField{Propósito}{Explicar máximo componente a componente e incremento propio.}
\FrameField{Qué decir}{Realice las tres ecuaciones. Destaque que el componente del líder cuenta eventos locales, no el valor del reloj Lamport.}
\FrameField{Pregunta}{¿Por qué se incrementa después del máximo?}
\FrameField{Respuesta esperada}{Porque recibir es un nuevo evento local posterior a ambos pasados.}
\FrameField{Cierre/transición}{Ahora interpretaremos los vectores en la historia del negocio.}
\FrameField{Tiempo sugerido}{3 min}

\subsection*{C08-C07 --- Vector clocks en la historia}
\FrameField{Propósito}{Conectar metadatos con decisiones de versión.}
\FrameField{Qué decir}{Explique dominación, igualdad e incomparabilidad. Anticipe que Dynamo usará la genealogía para no descartar ramas concurrentes.}
\FrameField{Pregunta}{¿Detectar concurrencia resuelve el conflicto?}
\FrameField{Respuesta esperada}{No; solo informa que no hay posterioridad causal. El dominio necesita merge, autoridad o rechazo.}
\FrameField{Cierre/transición}{Cerramos el bloque comparando los tres tipos de reloj.}
\FrameField{Tiempo sugerido}{2 min}

\subsection*{C08-C08 --- Tres relojes, tres preguntas}
\FrameField{Propósito}{Evitar confusión entre tiempo civil, duración y orden lógico.}
\FrameField{Qué decir}{Defina wall clock para fecha, monotonic para elapsed time y logical clock para orden. Use timeout como ejemplo del monotónico.}
\FrameField{Pregunta}{¿Cuál usaría para medir 500 ms de espera?}
\FrameField{Respuesta esperada}{Monotónico. ¿Para causalidad? Lamport/vector.}
\FrameField{Cierre/transición}{Los relojes describen; ahora necesitamos decidir una transición.}
\FrameField{Tiempo sugerido}{2 min}

\section{Bloque: Máquinas de estado y consenso}

\subsection*{C08-S01 --- Máquina de estados: una regla determinista}
\FrameField{Propósito}{Recordar la abstracción antes de replicarla.}
\FrameField{Qué decir}{Dibuje AVAILABLE y las dos reservas. Explique que igual estado+igual comando produce igual nuevo estado/respuesta.}
\FrameField{Pregunta}{¿Qué responde Reserve(Tomás) después de \texttt{RESERVED\_BY\_SOFIA}?}
\FrameField{Respuesta esperada}{Rechazo sin cambiar propietario.}
\FrameField{Cierre/transición}{La replicación exige que todas apliquen el mismo orden.}
\FrameField{Tiempo sugerido}{2 min}

\subsection*{C08-S02 --- Máquina de estados replicada}
\FrameField{Propósito}{Mostrar que cada nodo contiene log, código y estado local.}
\FrameField{Qué decir}{No trate la máquina como actor externo. Cada réplica aplica entradas committed en orden. La convergencia depende de determinismo y secuencia común.}
\FrameField{Pregunta}{¿Qué pasa si N2 aplica Sofía→Tomás y N3 Tomás→Sofía?}
\FrameField{Respuesta esperada}{Pueden terminar con estados/respuestas diferentes.}
\FrameField{Cierre/transición}{Formalizaremos dos familias de promesas: safety y liveness.}
\FrameField{Tiempo sugerido}{2 min}

\subsection*{C08-S03 --- Safety y liveness}
\FrameField{Propósito}{Definir propiedades sin vaguedad.}
\FrameField{Qué decir}{Safety: nada incorrecto; liveness: eventualmente progresa bajo condiciones. Use VIP-01 y mayoría estable como ejemplos.}
\FrameField{Pregunta}{¿Puede conservarse safety sin liveness?}
\FrameField{Respuesta esperada}{Sí: una minoría rechaza/espera indefinidamente y no produce doble venta.}
\FrameField{Cierre/transición}{Raft organiza liderazgo y log para sostener estas promesas.}
\FrameField{Tiempo sugerido}{2 min}

\subsection*{C08-S04 --- Raft: término, índice y entrada de log}
\FrameField{Propósito}{Definir vocabulario justo a tiempo.}
\FrameField{Qué decir}{Term es ronda creciente de liderazgo; index es posición; entry incluye term+command. Explique que 12 y 87 son valores didácticos, no horas.}
\FrameField{Pregunta}{¿Por qué term e index son distintos?}
\FrameField{Respuesta esperada}{Un término puede crear muchas entradas; el índice identifica la posición global del log.}
\FrameField{Cierre/transición}{Los términos permiten reconocer mensajes antiguos.}
\FrameField{Tiempo sugerido}{2 min}

\subsection*{C08-S05 --- Mensaje o líder obsoleto}
\FrameField{Propósito}{Explicar autoridad temporal y rechazo.}
\FrameField{Qué decir}{Muestre N1 en term 12 despertando cuando el cluster ya conoce 13. Un receptor actualiza/rechaza según término; no confíe en que el viejo líder “se dé cuenta” solo.}
\FrameField{Pregunta}{¿Qué evita que N1 siga mandando dentro del cluster?}
\FrameField{Respuesta esperada}{Las reglas de término; para recursos externos puede requerirse fencing.}
\FrameField{Cierre/transición}{Comenzamos la secuencia de commit con una sola copia.}
\FrameField{Tiempo sugerido}{2 min}

\subsection*{C08-S06 --- Raft paso 1: la propuesta existe solo en N1}
\FrameField{Propósito}{Distinguir append local de decisión.}
\FrameField{Qué decir}{N1 recibe Reserve(Sofía), crea entry 87/12 y la persiste localmente. Aún no puede decir committed ni responder confirmada.}
\FrameField{Pregunta}{¿Qué ocurre si N1 cae ahora?}
\FrameField{Respuesta esperada}{La entrada puede no sobrevivir a la elección; no había mayoría.}
\FrameField{Cierre/transición}{Agreguemos una segunda copia.}
\FrameField{Tiempo sugerido}{2 min}

\subsection*{C08-S07 --- Raft paso 2: dos copias todavía no son mayoría}
\FrameField{Propósito}{Explicar intersección de mayorías en cinco nodos.}
\FrameField{Qué decir}{N1 y N2 son 2/5. Los otros tres aún podrían formar una mayoría futura sin conocer la entrada. Evite decir solo “más copias es mejor”.}
\FrameField{Pregunta}{¿Qué conjunto futuro puede excluir a N1 y N2?}
\FrameField{Respuesta esperada}{N3,N4,N5.}
\FrameField{Cierre/transición}{La tercera réplica cambia esa geometría.}
\FrameField{Tiempo sugerido}{2 min}

\subsection*{C08-S08 --- Raft paso 3: la tercera réplica cambia la situación}
\FrameField{Propósito}{Justificar 3/5 y no solo memorizarlo.}
\FrameField{Qué decir}{Toda mayoría de tres intersecta con {N1,N2,N3}. Bajo las reglas de Raft, la entrada puede llegar a commit cuando el líder sabe que una mayoría la almacenó.}
\FrameField{Pregunta}{¿La mayoría basta sin reglas de términos/log matching?}
\FrameField{Respuesta esperada}{No; es ingrediente de la prueba, no garantía aislada.}
\FrameField{Cierre/transición}{Ahora separemos commit, apply y ACK.}
\FrameField{Tiempo sugerido}{2 min}

\subsection*{C08-S09 --- Raft paso 4: quién compromete, aplica y responde}
\FrameField{Propósito}{Asignar cada acción a un nodo.}
\FrameField{Qué decir}{N1 detecta mayoría, avanza commitIndex, aplica en su máquina local y responde. Followers aplican al conocer leaderCommit. Recorra las flechas en ese orden.}
\FrameField{Pregunta}{¿Puede N2 aplicar antes de saber que está committed?}
\FrameField{Respuesta esperada}{No debe presentar la entrada como decisión aplicada autoritativa solo por haberla replicado.}
\FrameField{Cierre/transición}{Moveremos el crash para probar la frontera.}
\FrameField{Tiempo sugerido}{3 min}

\subsection*{C08-S10 --- Crash antes y después de la mayoría}
\FrameField{Propósito}{Usar contrafactuales para verificar seguridad.}
\FrameField{Qué decir}{Compare crash con 1/5 y con 3/5. En el segundo caso diferencie ACK perdido de decisión perdida; el retry debe ser idempotente.}
\FrameField{Pregunta}{¿Qué debe devolver un nuevo líder ante el retry de Sofía?}
\FrameField{Respuesta esperada}{El resultado consistente con la entrada ya committed.}
\FrameField{Cierre/transición}{Una secuencia de commits construye un servicio de orden total.}
\FrameField{Tiempo sugerido}{3 min}

\subsection*{C08-S11 --- Total-order broadcast: por qué aparece aquí}
\FrameField{Propósito}{Justificar la abstracción desde máquinas de estados.}
\FrameField{Qué decir}{Diga “mismos mensajes committed, mismo orden”. Consenso por cada posición produce el log. No presente atomic broadcast como producto nuevo.}
\FrameField{Pregunta}{¿Qué falta si solo tenemos timestamps Lamport?}
\FrameField{Respuesta esperada}{Entrega común, selección de conjunto y ausencia de huecos.}
\FrameField{Cierre/transición}{Las pausas revelan por qué autoridad necesita números crecientes.}
\FrameField{Tiempo sugerido}{2 min}

\subsection*{C08-S12 --- Pausa de proceso, GC y fencing}
\FrameField{Propósito}{Definir GC y proteger recursos externos.}
\FrameField{Qué decir}{Narree N1 pausado, N2 elegido en term 13 y N1 despertando. Dentro de Raft, term obsoleto; fuera, token creciente que el recurso recuerda.}
\FrameField{Pregunta}{¿Por qué un lease basado solo en reloj local es riesgoso?}
\FrameField{Respuesta esperada}{Pausas y relojes pueden hacer que un proceso actúe después de perder autoridad.}
\FrameField{Cierre/transición}{Paxos permite transferir la misma meta a otro vocabulario.}
\FrameField{Tiempo sugerido}{3 min}

\subsection*{C08-S13 --- Paxos: misma meta, organización diferente}
\FrameField{Propósito}{Comparar sin enseñar dos algoritmos completos.}
\FrameField{Qué decir}{Use una tabla: term/ballot, index/slot, leader/proposer, follower/acceptor. Ambos impiden dos valores elegidos por posición; Raft organiza un log y liderazgo de modo distinto.}
\FrameField{Pregunta}{¿Por qué hablamos de Multi-Paxos?}
\FrameField{Respuesta esperada}{Porque un sistema necesita decidir muchos slots y reutilizar liderazgo para el camino normal.}
\FrameField{Cierre/transición}{Veamos el mapa en Kubernetes/etcd.}
\FrameField{Tiempo sugerido}{2 min}

\subsection*{C08-S14 --- Caso real: dónde aparecen estas ideas en Kubernetes}
\FrameField{Propósito}{Cerrar consenso con una arquitectura reconocible.}
\FrameField{Qué decir}{Recorra kubectl→API server→etcd líder→followers→KV committed. Mapee log, mayoría, máquina de estados, WAL/snapshots.}
\FrameField{Pregunta}{¿Qué se pierde si etcd pierde mayoría?}
\FrameField{Respuesta esperada}{La capacidad normal de comprometer nuevos cambios del control plane, no necesariamente todos los pods existentes de inmediato.}
\FrameField{Cierre/transición}{Antes de CAP, separemos tres clases de log.}
\FrameField{Tiempo sugerido}{2 min}

\subsection*{C08-S15 --- Tres logs que no debemos confundir}
\FrameField{Propósito}{Consolidar fronteras entre clases 6, 7 y 8.}
\FrameField{Qué decir}{Consenso decide comandos; WAL permite recovery local; Kafka distribuye hechos. Dé un ejemplo de cada uno para la reserva.}
\FrameField{Pregunta}{¿Publicar OrderConfirmed vuelve committed la reserva?}
\FrameField{Respuesta esperada}{No; la publicación debe seguir una decisión autoritativa y segura.}
\FrameField{Cierre/transición}{Ahora partimos la red después del commit de Sofía.}
\FrameField{Tiempo sugerido}{2 min}

\section{Bloque: CAP, PACELC y garantías}

\subsection*{C08-P01 --- La partición ocurre después del commit de Sofía}
\FrameField{Propósito}{Construir el escenario exacto de CAP.}
\FrameField{Qué decir}{Marque mayoría 3 y minoría 2. Sofía ya está committed; Tomás llega a la minoría, que sigue sana localmente.}
\FrameField{Pregunta}{¿Por qué la minoría no puede simplemente leer su copia?}
\FrameField{Respuesta esperada}{Porque puede estar atrasada y no sabe si la otra partición comprometió algo.}
\FrameField{Cierre/transición}{Formalizamos esa falta de conocimiento con historias indistinguibles.}
\FrameField{Tiempo sugerido}{2 min}

\subsection*{C08-P02 --- Dos historias indistinguibles para la minoría}
\FrameField{Propósito}{Derivar CAP por información, no por slogan.}
\FrameField{Qué decir}{Compare H1 (Sofía committed) y H2 (no ocurrió). Para la minoría, los mensajes observados son iguales durante la partición.}
\FrameField{Pregunta}{¿Puede una respuesta local ser correcta en ambas historias?}
\FrameField{Respuesta esperada}{No si debe confirmar en H2 y preservar una única historia en H1.}
\FrameField{Cierre/transición}{La siguiente lámina nombra las dos políticas.}
\FrameField{Tiempo sugerido}{2 min}

\subsection*{C08-P03 --- CAP: la decisión durante una partición}
\FrameField{Propósito}{Nombrar C/A después de la derivación.}
\FrameField{Qué decir}{Responder ambos lados prioriza disponibilidad y arriesga contradicción; esperar/rechazar preserva semántica fuerte y sacrifica disponibilidad para esa operación.}
\FrameField{Pregunta}{¿CAP clasifica para siempre a todo el producto?}
\FrameField{Respuesta esperada}{No; la decisión es por operación y se activa durante partición.}
\FrameField{Cierre/transición}{Distingamos disponibilidad CAP de otras métricas.}
\FrameField{Tiempo sugerido}{2 min}

\subsection*{C08-P04 --- Uptime, disponibilidad CAP y disponibilidad semántica}
\FrameField{Propósito}{Evitar confundir health con éxito de operación.}
\FrameField{Qué decir}{Defina uptime por ventana, CAP como respuesta no errónea y semántica como cumplimiento de promesa de negocio. Use 503 como contraejemplo.}
\FrameField{Pregunta}{¿Por qué se habla de uptime anual?}
\FrameField{Respuesta esperada}{Es una ventana contractual común, no parte del teorema.}
\FrameField{Cierre/transición}{Incluso sin partición, la coordinación sigue costando.}
\FrameField{Tiempo sugerido}{2 min}

\subsection*{C08-P05 --- PACELC: el costo continúa cuando la red funciona}
\FrameField{Propósito}{Extender el análisis a operación normal.}
\FrameField{Qué decir}{Compare lectura local 15 ms y coordinada 180 ms. Pregunte qué operación necesita cuál. No trate PACELC como clasificación de producto.}
\FrameField{Pregunta}{¿Qué medir además de latencia?}
\FrameField{Respuesta esperada}{Antigüedad/staleness, disponibilidad, errores e impacto de coordinar.}
\FrameField{Cierre/transición}{Necesitamos vocabulario preciso de consistencia.}
\FrameField{Tiempo sugerido}{2 min}

\subsection*{C08-P06 --- Recencia, linearizabilidad y serializabilidad}
\FrameField{Propósito}{Distinguir garantías frecuentemente mezcladas.}
\FrameField{Qué decir}{Use historias mínimas: frescura, write terminado antes de read, transacciones equivalentes a orden serial, y orden serial compatible con tiempo real externo.}
\FrameField{Pregunta}{¿Serializabilidad implica que una lectura ve el último write real?}
\FrameField{Respuesta esperada}{No necesariamente; el orden serial puede no respetar el tiempo real externo.}
\FrameField{Cierre/transición}{Entramos ahora en quórums Dynamo-like.}
\FrameField{Tiempo sugerido}{3 min}

\section{Bloque: Quórums, versiones y Dynamo}

\subsection*{C08-D01 --- Mayoría en consenso y quorum Dynamo-like no son sinónimos}
\FrameField{Propósito}{Prevenir transferencia incorrecta de “3 respuestas”.}
\FrameField{Qué decir}{Recorra fila por fila: pregunta, unidad, reglas, concurrencia, resultado e intersección. En Dynamo pueden quedar varias versiones.}
\FrameField{Pregunta}{¿Qué entrega una lectura Dynamo si encuentra siblings?}
\FrameField{Respuesta esperada}{Una política debe reconciliarlos o exponerlos; no hay único commit por slot.}
\FrameField{Cierre/transición}{Primero motivaremos la necesidad combinatoria de intersección.}
\FrameField{Tiempo sugerido}{3 min}

\subsection*{C08-D02 --- Por qué necesitamos intersección}
\FrameField{Propósito}{Mostrar el problema antes de la fórmula.}
\FrameField{Qué decir}{Escriba en N1,N2 y lea N4,N5. Los conjuntos disjuntos permiten stale read aun con operaciones “exitosas”.}
\FrameField{Pregunta}{¿Qué propiedad buscamos al aumentar R o W?}
\FrameField{Respuesta esperada}{Que toda lectura ideal cruce al menos una réplica de toda escritura ideal.}
\FrameField{Cierre/transición}{Derivamos la desigualdad y sus límites.}
\FrameField{Tiempo sugerido}{2 min}

\subsection*{C08-D03 --- N, R y W: construir la intersección}
\FrameField{Propósito}{Definir N,R,W y visualizar conjuntos.}
\FrameField{Qué decir}{Use N=5 y pinte conjuntos. Explique disponibilidad de lectura/escritura para distintas configuraciones. No prometa linearizabilidad.}
\FrameField{Pregunta}{¿Qué configuración favorece lecturas y cuál escrituras?}
\FrameField{Respuesta esperada}{R pequeño favorece lectura; W pequeño escritura, con otros riesgos.}
\FrameField{Cierre/transición}{Tres contraejemplos mostrarán por qué la fórmula no basta.}
\FrameField{Tiempo sugerido}{3 min}

\subsection*{C08-D04 --- Contraejemplo 1: escrituras concurrentes}
\FrameField{Propósito}{Mostrar que el nodo común puede conservar conflicto.}
\FrameField{Qué decir}{Dos coordinadores escriben ramas concurrentes. La intersección ayuda a descubrir ambas, no a decidir cuál es posterior.}
\FrameField{Pregunta}{¿Qué metadato permite detectar la concurrencia?}
\FrameField{Respuesta esperada}{Versiones/vector clocks u otro esquema causal.}
\FrameField{Cierre/transición}{El siguiente contraejemplo cambia el conjunto efectivo.}
\FrameField{Tiempo sugerido}{2 min}

\subsection*{C08-D05 --- Contraejemplo 2: sloppy quorum}
\FrameField{Propósito}{Mostrar que el quorum real puede usar nodos alternativos.}
\FrameField{Qué decir}{Con propietarios inaccesibles, la escritura se acepta fuera de la preference list. Antes del handoff, una lectura sobre propietarios puede no intersectar efectivamente.}
\FrameField{Pregunta}{¿Qué deuda crea esta disponibilidad?}
\FrameField{Respuesta esperada}{Hints, datos fuera de lugar, frescura variable y repair.}
\FrameField{Cierre/transición}{El tercer caso muestra el peligro de elegir por reloj físico.}
\FrameField{Tiempo sugerido}{2 min}

\subsection*{C08-D06 --- Contraejemplo 3: last-write-wins y relojes desalineados}
\FrameField{Propósito}{Mostrar pérdida de información por timestamps.}
\FrameField{Qué decir}{Construya dos edits donde la posterior causalmente tiene timestamp menor. LWW descarta la actualización válida si confía en reloj desalineado.}
\FrameField{Pregunta}{¿Cuándo LWW podría ser aceptable?}
\FrameField{Respuesta esperada}{Solo si el dominio acepta pérdida y la regla está explícita; no como garantía causal general.}
\FrameField{Cierre/transición}{Volvemos a vector clocks, ahora sobre versiones.}
\FrameField{Tiempo sugerido}{2 min}

\subsection*{C08-D07 --- Vector clocks reaparecen como genealogía del carrito}
\FrameField{Propósito}{Reutilizar causalidad en un contexto de datos.}
\FrameField{Qué decir}{Parta de v0 y cree ramas [1,0] y [0,1]. Explique que son siblings concurrentes y que una lectura/reconciliación puede producir una versión que domine ambas.}
\FrameField{Pregunta}{¿Qué diferencia hay con los vectores del timeline?}
\FrameField{Respuesta esperada}{La matemática es la misma; ahora resume genealogía de versiones persistidas.}
\FrameField{Cierre/transición}{El dominio decide si el merge es legal.}
\FrameField{Tiempo sugerido}{2 min}

\subsection*{C08-D08 --- Carrito fusionable; entrada VIP no fusionable}
\FrameField{Propósito}{Enseñar que consistencia es semántica.}
\FrameField{Qué decir}{Unión de ítems puede conservar intención; dos propietarios no. Mencione cantidades, removals y promociones como límites de un merge ingenuo.}
\FrameField{Pregunta}{¿Detectar conflicto es suficiente?}
\FrameField{Respuesta esperada}{No; necesita regla de dominio, autoridad o intervención.}
\FrameField{Cierre/transición}{Ahora recorremos el sistema Dynamo completo.}
\FrameField{Tiempo sugerido}{2 min}

\subsection*{C08-D09 --- Dynamo: localizar, escribir, conservar y reparar}
\FrameField{Propósito}{Presentar la arquitectura como flujo causal.}
\FrameField{Qué decir}{Defina hash, preference list, N, coordinador, W y versiones. Recorra una escritura normal antes de fallas.}
\FrameField{Pregunta}{¿Qué parte distribuye carga y cuál maneja conflicto?}
\FrameField{Respuesta esperada}{Hash/preference list distribuyen; versioning/reconciliation manejan conflicto.}
\FrameField{Cierre/transición}{Luego falla un propietario y aparece sloppy quorum.}
\FrameField{Tiempo sugerido}{3 min}

\subsection*{C08-D10 --- Sloppy quorum, hint y hinted handoff}
\FrameField{Propósito}{Definir cada término mediante una ejecución.}
\FrameField{Qué decir}{El alternativo acepta la escritura y guarda “esto pertenece a B”. Cuando B vuelve, transfiere. Aclare que hint no es copia autoritativa eterna.}
\FrameField{Pregunta}{¿Qué mediría durante una falla larga?}
\FrameField{Respuesta esperada}{Cantidad/bytes/edad de hints y tasa de handoff.}
\FrameField{Cierre/transición}{Otras divergencias requieren repair durante lectura o background.}
\FrameField{Tiempo sugerido}{2 min}

\subsection*{C08-D11 --- Read repair, anti-entropy y tasa de divergencia}
\FrameField{Propósito}{Distinguir tres mecanismos y la deuda.}
\FrameField{Qué decir}{Read repair corrige lo leído; anti-entropy compara rangos; handoff devuelve desplazados. Defina divergencia nueva y capacidad resuelta.}
\FrameField{Pregunta}{¿Qué dato puede quedar sin reparar si nadie lo lee?}
\FrameField{Respuesta esperada}{Dato frío; por eso se necesita anti-entropy.}
\FrameField{Cierre/transición}{Cerramos Dynamo explicando por qué aceptó esos costos.}
\FrameField{Tiempo sugerido}{3 min}

\subsection*{C08-D12 --- Por qué Dynamo aceptó esa complejidad}
\FrameField{Propósito}{Conectar mecanismos con prioridad always-on.}
\FrameField{Qué decir}{Diga que no es complejidad gratuita: cada mecanismo sostiene disponibilidad de servicios concretos. Enumere conflictos, hints, espacio, repair y semántica como costos.}
\FrameField{Pregunta}{¿Para qué operación de nuestra clase sería inadecuado?}
\FrameField{Respuesta esperada}{VIP-01 si la política permite dos confirmaciones y merge posterior.}
\FrameField{Cierre/transición}{Spanner aparece cuando necesitamos transacciones globales fuertes.}
\FrameField{Tiempo sugerido}{2 min}

\section{Bloque: Spanner, TrueTime y orden externo}

\subsection*{C08-G01 --- Por qué aparece Spanner en la misma historia}
\FrameField{Propósito}{Reanclar Spanner en la compra.}
\FrameField{Qué decir}{La compra toca entradas, pagos y derechos en grupos distintos; además, una acción posterior en Madrid debe respetar la confirmación vista en BA.}
\FrameField{Pregunta}{¿Qué dos problemas nuevos aparecen?}
\FrameField{Respuesta esperada}{Atomicidad multi-grupo y orden externo compatible con observación del cliente.}
\FrameField{Cierre/transición}{Primero se coordina por grupos, no con consenso global único.}
\FrameField{Tiempo sugerido}{2 min}

\subsection*{C08-G02 --- Spanner: consenso localizado y transacción multi-grupo}
\FrameField{Propósito}{Explicar partición de coordinación.}
\FrameField{Qué decir}{Cada rango/grupo tiene replicas y consenso. Una transacción coordina solo grupos participantes, mediante protocolo transaccional sobre esos líderes.}
\FrameField{Pregunta}{¿Por qué no una votación planetaria por operación?}
\FrameField{Respuesta esperada}{Escalaría mal y coordinaría datos irrelevantes; la autoridad se localiza por rango.}
\FrameField{Cierre/transición}{Todavía falta ordenar commits respecto del tiempo observado.}
\FrameField{Tiempo sugerido}{2 min}

\subsection*{C08-G03 --- TrueTime: intervalo de tiempo con incertidumbre explícita}
\FrameField{Propósito}{Definir API e infraestructura de sincronización.}
\FrameField{Qué decir}{Diga que GPS/relojes atómicos alimentan time masters y servidores sincronizan. earliest/latest acotan la hora real; no son dos relojes de negocio.}
\FrameField{Pregunta}{¿Qué hace crecer epsilon?}
\FrameField{Respuesta esperada}{Error, deriva y tiempo desde última sincronización.}
\FrameField{Cierre/transición}{El coordinador usa esa cota para elegir un timestamp.}
\FrameField{Tiempo sugerido}{3 min}

\subsection*{C08-G04 --- Quién asigna el timestamp de commit}
\FrameField{Propósito}{Asignar responsabilidad y restricciones.}
\FrameField{Qué decir}{El coordinador elige s respetando prepares y TT.latest. No diga que TrueTime “asigna” el commit; proporciona una cota.}
\FrameField{Pregunta}{¿Por qué elegir s dentro o por encima de latest?}
\FrameField{Respuesta esperada}{Para ubicar el commit de forma segura respecto del intervalo y participantes.}
\FrameField{Cierre/transición}{Después se paga commit wait.}
\FrameField{Tiempo sugerido}{2 min}

\subsection*{C08-G05 --- Commit wait: esperar hasta demostrar que $s$ quedó en el pasado}
\FrameField{Propósito}{Explicar la desigualdad con números.}
\FrameField{Qué decir}{Con s=16 y TT=[4,16], aún no sabemos que 16 pasó. Espere hasta earliest>16, por ejemplo [17,29]. Mayor incertidumbre puede alargar la espera.}
\FrameField{Pregunta}{¿Qué compra esa espera?}
\FrameField{Respuesta esperada}{Que cuando el commit se hace visible, su timestamp ya está en el pasado real.}
\FrameField{Cierre/transición}{Esto sostiene una consistencia que respeta lo visto fuera del motor.}
\FrameField{Tiempo sugerido}{3 min}

\subsection*{C08-G06 --- Por qué se llama consistencia externa}
\FrameField{Propósito}{Definir “externa” con dos transacciones.}
\FrameField{Qué decir}{Si el cliente recibe T1 y luego inicia T2, la serialización debe respetar T1<T2. No basta cualquier orden serial interno.}
\FrameField{Pregunta}{¿En qué se diferencia de serializabilidad sola?}
\FrameField{Respuesta esperada}{Agrega compatibilidad con el orden real observado externamente.}
\FrameField{Cierre/transición}{Cerramos mapeando cada mecanismo a su pregunta.}
\FrameField{Tiempo sugerido}{2 min}

\section{Bloque: Integración y cierre}

\subsection*{C08-I01 --- Cada mecanismo responde una pregunta distinta}
\FrameField{Propósito}{Sintetizar sin reducir a slogans.}
\FrameField{Qué decir}{Recorra la tabla y pida un ejemplo de confusión: Kafka vs consenso, vector vs decisión, WAL vs replicación.}
\FrameField{Pregunta}{¿Qué mecanismo usarían para detectar concurrencia? ¿Para decidir autoridad?}
\FrameField{Respuesta esperada}{Vectores; consenso.}
\FrameField{Cierre/transición}{La actividad final exige combinar decisiones por operación.}
\FrameField{Tiempo sugerido}{3 min}

\subsection*{C08-I02 --- Actividad integradora: cuatro operaciones, cuatro decisiones}
\FrameField{Propósito}{Transferir el criterio a otros dominios.}
\FrameField{Qué decir}{Pida grupos de dos. Para cada operación deben declarar invariante, mecanismo, respuesta bajo partición, costo normal y prueba.}
\FrameField{Pregunta}{¿Qué operación admite merge y cuál requiere autoridad?}
\FrameField{Respuesta esperada}{Carrito puede; entrada/ledger/premio suelen requerir autoridad fuerte.}
\FrameField{Cierre/transición}{El exit ticket recoge evidencia individual.}
\FrameField{Tiempo sugerido}{8 min}

\subsection*{C08-I03 --- Exit ticket}
\FrameField{Propósito}{Evaluar comprensión esencial.}
\FrameField{Qué decir}{Dé tres minutos individuales. No corrija cada frase oralmente; recoja y use los errores para abrir la Clase 09.}
\FrameField{Pregunta}{¿Qué respuesta distinguiría mejor comprensión superficial de profunda?}
\FrameField{Respuesta esperada}{La diferencia R+W>N vs linearizabilidad o log de consenso vs Kafka.}
\FrameField{Cierre/transición}{Cierre con la bibliografía como ruta, no como lista decorativa.}
\FrameField{Tiempo sugerido}{3 min}

\subsection*{C08-I04 --- Bibliografía guiada}
\FrameField{Propósito}{Orientar estudio posterior.}
\FrameField{Qué decir}{Explique qué aporta cada fuente: DDIA integra; Lamport causalidad/Paxos; Raft comprensibilidad; Dynamo y Spanner decisiones de producción.}
\FrameField{Pregunta}{¿Qué lectura conviene para la duda que quedó?}
\FrameField{Respuesta esperada}{Elija según causalidad, consenso, disponibilidad o tiempo global.}
\FrameField{Cierre/transición}{Cerrar recordando la pregunta transversal de costo del error.}
\FrameField{Tiempo sugerido}{1 min}

\section{Demostraciones guiadas}
\subsection{Demo 1: causalidad y relojes}
Tres estudiantes representan BA, líder y Madrid. Entregue tarjetas de mensaje y haga actualizar Lamport y vector en voz alta. Error a buscar: copiar el 5 escalar dentro del vector.

\subsection{Demo 2: Raft humano}
Cinco estudiantes sostienen logs. El líder entrega entry 87 a N2 y N3. Mueva CRASH antes y después de la mayoría. Pida justificar cada ACK.

\subsection{Demo 3: quorum y versiones}
Coloree conjuntos R y W. Después entregue dos versiones concurrentes del carrito y pida un merge; repita con VIP-01 para mostrar el límite semántico.

\subsection{Demo 4: TrueTime}
Muestre intervalos sucesivos y un timestamp \(s\). Los alumnos levantan tarjeta ESPERAR hasta que \(earliest>s\).

\section{Preguntas socráticas}
\begin{longtable}{p{0.43\textwidth}p{0.49\textwidth}}
\toprule
Pregunta & Núcleo esperado\\
\midrule
¿Qué sabe realmente un nodo tras timeout? & Solo que no recibió respuesta en su plazo.\\
¿Por qué un vector incomparables importa? & Revela ramas concurrentes; no elige ganador.\\
¿Cuándo una entrada es committed? & Cuando el protocolo satisface la regla de mayoría y seguridad, no solo por existir copias.\\
¿Puede haber safety sin liveness? & Sí; una minoría aislada puede no progresar y no contradecir.\\
¿Qué hace CAP visible? & La imposibilidad de responder en dos historias indistinguibles sin romper una garantía.\\
¿Qué garantiza R+W>N? & Intersección ideal, no orden ni linearizabilidad completa.\\
¿Qué deuda compra sloppy quorum? & Hints, datos fuera de lugar, divergencia y repair.\\
¿Por qué TrueTime usa intervalo? & Para declarar incertidumbre en lugar de ocultarla.\\
\bottomrule
\end{longtable}

\section{Banco de preguntas previsibles}
\begin{description}
\item[¿Por qué no usar solo el reloj del servidor?] Porque sincronización y pausas tienen error; un timestamp físico no prueba causalidad ni orden autoritativo.
\item[¿Raft siempre necesita disco?] Su safety presupone que el estado persistente especificado sobreviva reinicios; la implementación concreta puede usar distintos medios, pero no puede olvidar votos/entradas que prometió recordar.
\item[¿Kafka usa Raft?] Las versiones modernas usan un protocolo basado en Raft para metadatos del cluster, pero el repaso de la clase trata principalmente el log de datos y sus réplicas. No confundir plano de control con semántica de negocio.
\item[¿CAP ocurre todo el tiempo?] El trade-off específico se activa durante una partición; PACELC recuerda que coordinar también cuesta con red sana.
\item[¿Dynamo es DynamoDB?] No. La clase usa el diseño del paper original como caso histórico y conceptual.
\item[¿Spanner viola CAP?] No. Durante particiones puede sacrificar disponibilidad de operaciones que requieren coordinación fuerte.
\end{description}

\section{Errores previsibles}
\begin{longtable}{p{0.38\textwidth}p{0.54\textwidth}}
\toprule
Error & Intervención\\
\midrule
“Timeout significa caída.” & Pedir cuatro historias compatibles.\\
“Lamport menor significa causa.” & Construir dos procesos sin mensajes.\\
“El vector del líder termina en 5.” & Separar etiqueta escalar de conteo por actor.\\
“Tres copias siempre implican commit.” & Preguntar protocolo, término y conjunto futuro.\\
“CAP es elegir dos de tres.” & Volver a H1/H2 durante partición concreta.\\
“503 rápido es disponibilidad.” & Separar uptime de respuesta válida CAP.\\
“R+W>N implica linearizabilidad.” & Usar escrituras concurrentes o sloppy quorum.\\
“Eventual significa que se arregla solo.” & Exigir tasa de repair, edad y SLO de convergencia.\\
“TrueTime da la hora exacta.” & Dibujar intervalo y hacer crecer epsilon.\\
\bottomrule
\end{longtable}

\section{Contingencias}
\begin{description}
\item[El grupo no participa.] Use voto A/B/C, 30 segundos individuales y comparación en parejas antes de pedir una explicación pública.
\item[Se consume demasiado tiempo en Kafka.] Cierre con una frase: ``Kafka preserva registros; ahora necesitamos decidir una única transición autoritativa'' y salte a VIP-01.
\item[Lamport y vectores se confunden.] Mantenga simultáneamente el escalar $L=5$ y el vector $[2,2,3]$; pregunte qué cuenta cada número.
\item[Raft resulta abstracto.] Use cinco tarjetas físicas y mueva CRASH antes y después del tercer ACK. No avance hasta separar replicated, committed y applied.
\item[CAP se reduce a dos de tres.] Vuelva a las dos historias indistinguibles para la minoría y pida una respuesta concreta a Tomás.
\item[No alcanza el tiempo.] Preserve causalidad, Raft, CAP, un contraejemplo de quórums, Dynamo como deuda y TrueTime. Use Paxos detallado y anti-entropy como respaldo.
\end{description}

\section{Parte C: auditoría}
La auditoría pedagógica debe comprobar que el caso VIP-01 permanece visible, que toda definición aparece después de una necesidad y que las actividades producen evidencia observable. Conviene conservar exit tickets, fotografías de timelines y respuestas a la actividad integradora.

\section{Preguntas que puede formular un auditor pedagógico}
\begin{itemize}
\item ¿Dónde se recupera una idea previa y cómo se verifica que fue comprendida?
\item ¿Qué producción observable reemplaza la simple exposición de una definición?
\item ¿Por qué cada nuevo término aparece en ese punto de la historia?
\item ¿Cómo se atiende a estudiantes sin bases sólidas de redes o sistemas operativos?
\item ¿Qué hace el docente cuando una predicción inicial es incorrecta?
\item ¿Cómo se distingue un caso real de una analogía y de una formalización?
\item ¿Qué slides pueden omitirse sin romper el hilo de tres horas?
\end{itemize}
La respuesta debe señalar diagnóstico, timelines incrementales, actividades de bajo riesgo, vocabulario justo a tiempo, libreto de transición y exit ticket.

\section{Criterios de evaluación formativa}
Una respuesta lograda incluye: operación, invariante, información local, falla, garantía, mecanismo, costo y evidencia. Niveles:
\begin{itemize}
\item \textbf{Inicial}: nombra tecnología sin ejecución.
\item \textbf{En desarrollo}: conecta mecanismo y garantía, pero no explicita supuestos/costo.
\item \textbf{Logrado}: reconstruye la ejecución y declara límite y métrica.
\item \textbf{Avanzado}: compara alternativa, construye contraejemplo y define condición de revisión.
\end{itemize}

\section{Parte B: dossier técnico}
Esta parte reúne las precisiones que el docente necesita para responder preguntas sin recargar las diapositivas: límites de relojes lógicos, persistencia de Raft, relación con Paxos, semántica CAP, contraejemplos de quórums, reparación en Dynamo y condiciones de TrueTime.

\section{Dossier técnico de preparación del docente}
\subsection{Cómo explicar los conceptos en tres profundidades}
\textbf{30 segundos:} causalidad describe relaciones; consenso elige una historia; CAP obliga a decidir durante partición; Dynamo conserva y repara; Spanner coordina y espera para orden externo.

\textbf{Dos minutos:} use VIP-01 y recorra click, mensajes concurrentes, entry 87, mayoría, partición 3|2, carrito concurrente y commit wait.

\textbf{Cinco minutos:} agregue límites: Lamport no detecta concurrencia, mayoría no es suficiente sin protocolo, R+W>N no implica linearizabilidad, merge depende del dominio y TrueTime paga incertidumbre con espera.

\section{Matriz de corrección y preguntas de defensa}
\begin{longtable}{p{0.23\textwidth}p{0.32\textwidth}p{0.37\textwidth}}
\toprule
Dimensión & Evidencia esperada & Pregunta de defensa\\
\midrule
Causalidad & timeline con flechas y concurrencia & ¿qué mensaje crea esta relación?\\
Autoridad & commit/owner claramente identificado & ¿quién puede confirmar y por qué?\\
Falla & crash/partición en punto exacto & ¿qué cambia si la falla se mueve?\\
Garantía & observación/invariante concreto & ¿qué historia está prohibida?\\
Costo & latencia, rechazo, repair o metadatos & ¿dónde se paga durante operación normal?\\
Métrica & p99, staleness, conflicts, repair lag & ¿qué valor obligaría a cambiar diseño?\\
\bottomrule
\end{longtable}

\section{Checklist personal de estudio}
Antes de dictar la clase, verifique que puede explicar sin leer:
\begin{enumerate}
\item cuatro historias compatibles con un timeout;
\item por qué $L=5$ y $[2,2,3]$ no se contradicen;
\item cuándo entry 87 pasa de local a committed y applied;
\item por qué una mayoría Raft no es el mismo concepto que $R$, $W$ y $N$;
\item un contraejemplo a ``$R+W>N$ implica linearizabilidad'';
\item la diferencia entre hinted handoff, read repair y anti-entropy;
\item quién elige $s$ y por qué se espera hasta $earliest>s$;
\item una transición oral clara entre cada bloque.
\end{enumerate}
Marque en la guía las slides de respaldo que omitirá si el ritmo se retrasa y prepare una respuesta breve para cada pregunta del banco.

\section{Trabajo Final Integrador}
Para cada subsistema del proyecto, pedir una ficha con \(D=(P,F,G,S,T)\): problema, fallas, garantías, solución y trade-offs. Debe identificar owner, clave/partición, frontera de commit, política bajo partición, métrica y prueba adversarial. No aceptar “usaremos consenso” o “será eventual” sin unidad, protocolo y condición de éxito.

\end{document}
